\begin{table}[t]
\centering
\begin{threeparttable}
\caption{Width of three interval constructions for the same ROC-AUC point estimates. The repeated-split band reported previously is 1.5--2.5x narrower than either sampling-uncertainty interval.}
\label{tab:tab-interval-methods}
\small
\begin{tabular}{lrrrrr}
\toprule
contrast & v1 split band width & Hanley 95\% width & bootstrap 95\% width & Hanley / split & bootstrap / split \\
\midrule
broad breach vs none & 0.104 & 0.185 & 0.237 & 1.77x & 2.28x \\
substantive leak vs rest & 0.133 & 0.174 & 0.181 & 1.31x & 1.36x \\
substantive leak vs appropriate & 0.152 & 0.185 & 0.233 & 1.22x & 1.53x \\
limiting vs direct (all) & 0.100 & 0.148 & 0.154 & 1.47x & 1.53x \\
limiting vs direct (disclosers only) & 0.094 & 0.153 & 0.168 & 1.63x & 1.78x \\
broad-only vs substantive & 0.108 & 0.181 & 0.195 & 1.68x & 1.81x \\
\bottomrule
\end{tabular}
\begin{tablenotes}[flushleft]\footnotesize
\item Widths are for ROC-AUC in the 216 analysis population.
\item The v1 band is the 2.5/97.5 percentile over 20 repeated cross-validation splits of the SAME cases. It measures split variability, not sampling uncertainty, and must never be presented as a confidence interval (Nadeau \& Bengio 2003; Bates, Hastie \& Tibshirani 2023).
\item Hanley \& McNeil (1982) assumes a single fixed scoring rule applied to independent cases; an average of per-fold AUCs from models refit on each training split violates that, so it is reported as an approximation alongside the bootstrap.
\item The bootstrap is a stratified prediction-resampling interval (n\_boot=1000), conditional on the fitted models; it is not a test against the null.
\end{tablenotes}
\end{threeparttable}
\end{table}
